\documentclass[12pt]{article}
\usepackage[utf8]{inputenc}
\usepackage[T1]{fontenc}
\usepackage[margin=1in]{geometry}
\usepackage{setspace}
\usepackage{xcolor}
\usepackage[colorlinks=true, urlcolor=blue!60!black, linkcolor=blue!60!black]{hyperref}
\setlength{\parindent}{0em}
\setlength{\parskip}{0.6em}
\renewcommand{\baselinestretch}{1.0}
\begin{document}
\begin{center}
{\LARGE \textit{Biden v.\ Nebraska}, 600 U.S.\ 477 (2023)\par}
\vspace{0.5em}
{\large Case Brief --- Two-Page Edition\par}
\end{center}

\section{Memory Jogger}
Missouri had standing through its loan servicer MOHELA, and the HEROES Act's power to ``waive or modify'' did not authorize \$430 billion in student-debt cancellation (\href{https://www.law.cornell.edu/supremecourt/text/22-506}{600 U.S. 477}).

\section{Detailed Case Facts}
The HEROES Act lets the Secretary of Education waive or modify Title IV provisions in a national emergency so borrowers are not left worse off. In 2022, invoking COVID--19, Secretary Cardona announced cancellation of up to \$10,000 (\$20,000 for Pell recipients) for borrowers earning under \$125,000, at an estimated \$430 billion cost. Six States sued; Missouri relied on lost fees to MOHELA, a state-created public instrumentality that did not join the suit.

\section{Procedural History}
The district court dismissed for lack of standing; the Eighth Circuit enjoined the plan; the Supreme Court granted certiorari before judgment and reversed.

\section{Judicial Votes}
6--3. Roberts wrote for Thomas, Alito, Gorsuch, Kavanaugh, and Barrett; Barrett also concurred; Kagan dissented, joined by Sotomayor and Jackson.

\section{Holding}
Financial injury to a state instrumentality that the State creates, controls, and may dissolve is injury to the State. ``Waive or modify'' permits only modest adjustments, not a new forgiveness regime, and major questions principles apply to benefits programs, requiring clear congressional authorization.

\section{Analysis of Opinions}
Roberts found MOHELA's \$44 million loss to be Missouri's own, relying on \textit{Arkansas v. Texas} and \textit{Lebron}. Following \textit{MCI}, ``modify'' means moderate change, and the Secretary waived nothing but instead wrote a new scheme. Under \textit{West Virginia v. EPA}, the plan's scale, cost, and political salience demanded clear authorization. Barrett recast the doctrine as a contextual, textualist tool, illustrated by her babysitter analogy. Kagan called the States ideological plaintiffs, argued that MOHELA's independence meant Missouri lost nothing, and read ``waive or modify'' as authorizing changes along a spectrum.

\section{Examples: Future Applications}
\textbf{Same-Side:} an agriculture secretary's \$350 billion farm-debt cancellation under a ``waive or modify'' clause fails. \textbf{Opposite-Side:} targeted hurricane waivers of deadlines and service rules, like the 2003 and 2012 waivers, are valid. \textbf{Fence-Sitter:} a three-year, \$150 billion payment and interest pause, a question Footnote 5 left open.

\section{Critique}
Critics argued that the MOHELA theory eroded Article III limits, a tension sharpened when the Tenth Circuit held MOHELA is not an arm of the State, and that the majority rewrote broad text. Empirical work by Tobia, Walters, and Slocum questioned Barrett's analogies, while conservatives such as Squitieri and Vermeule warned that contextualism slides into purposivism.

\section{Key Quotations}
\begin{quote}
``We hold today that the Act allows the Secretary to `waive or modify' existing statutory or regulatory provisions applicable to financial assistance programs under the Education Act, not to rewrite that statute from the ground up.''
\end{quote}
\noindent --- Chief Justice Roberts, Opinion of the Court, 600 U.S. 477.

\end{document}
